\documentclass[12pt,a4paper]{article}
\usepackage[utf8]{inputenc}
\usepackage[T1]{fontenc}
\usepackage[french]{babel}
\usepackage{geometry}
\usepackage{xcolor}
\usepackage{fancyhdr}

\geometry{margin=2cm}
\pagestyle{fancy}
\fancyhf{}
\fancyhead[L]{CEJM - Fiche de Révision}
\fancyhead[R]{Chapitre 1}
\fancyfoot[C]{\thepage}
\setlength{\headheight}{15pt}

% Couleurs
\definecolor{bleu}{RGB}{0,90,170}
\definecolor{vert}{RGB}{0,150,0}
\definecolor{rouge}{RGB}{200,0,0}
\definecolor{gris}{RGB}{240,240,240}

% Commandes pour encadrés
\newcommand{\definition}[2]{\noindent\fcolorbox{bleu}{gris}{\parbox{\dimexpr\textwidth-2\fboxsep-2\fboxrule\relax}{\textbf{\textcolor{bleu}{DÉFINITION :}} #1\\\textit{#2}}}}

\newcommand{\pointcle}[1]{\noindent\fcolorbox{vert}{white}{\parbox{\dimexpr\textwidth-2\fboxsep-2\fboxrule\relax}{\textbf{\textcolor{vert}{POINT CLÉ :}} #1}}}

\newcommand{\exemple}[1]{\noindent\fcolorbox{rouge}{white}{\parbox{\dimexpr\textwidth-2\fboxsep-2\fboxrule\relax}{\textbf{\textcolor{rouge}{EXEMPLE :}} #1}}}

\title{\textbf{Fiche de Révision CEJM}\\[2mm]\large Chapitre 1 : Les agents économiques en relation avec l'entreprise}
\author{}
\date{}

\begin{document}
\maketitle

\section*{Problématique}
\textbf{Comment s'établissent les relations entre l'entreprise et son environnement économique ?}

\section{Définitions essentielles}

\definition{Agent économique}{Individu ou ensemble d'individus constituant un centre de décision économique indépendant, regroupés selon leur fonction principale.}

\vspace{3mm}
\definition{Production marchande}{Biens et services vendus sur un marché à un prix concurrentiel, assurée par les entreprises.}

\vspace{3mm}
\definition{Production non marchande}{Biens et services fournis gratuitement ou à des prix non significatifs (< 50\% du coût), assurée par les administrations publiques.}

\vspace{3mm}
\definition{Circuit économique}{Modèle schématique des échanges et interdépendances entre agents économiques, représentant flux réels et monétaires.}

\vspace{3mm}
\definition{Marché des capitaux}{Lieu de rencontre entre offreurs (épargnants) et demandeurs de capitaux (entreprises), où le prix du capital est le taux d'intérêt.}

\section{Les 4 agents économiques principaux}

\subsection{Ménages}
\pointcle{Fonction principale : CONSOMMATION de biens et services pour satisfaire directement leurs besoins}
\begin{itemize}
    \item Individus ou groupes vivant dans un même logement
    \item Offrent leur travail sur le marché du travail
    \item Consomment et épargnent
\end{itemize}

\subsection{Entreprises}
\pointcle{Fonction principale : PRODUCTION de biens et/ou services pour les vendre et réaliser un profit}
\begin{itemize}
    \item Production marchande destinée à la vente
    \item Demandeurs de travail et de capitaux
    \item Relations B2B avec d'autres entreprises
\end{itemize}

\subsection{Banques}
\pointcle{Fonction principale : FINANCEMENT de l'économie par collecte et octroi de crédits}
\begin{itemize}
    \item Intermédiaires financiers
    \item Collectent l'épargne des ménages
    \item Accordent des crédits aux entreprises et ménages
\end{itemize}

\subsection{Administrations publiques}
\pointcle{Fonction principale : Répondre aux besoins d'intérêt général et produire des services non marchands}
\begin{itemize}
    \item État, Sécurité sociale, collectivités locales
    \item Fournissent services publics (éducation, sécurité, justice)
    \item Financés par impôts, taxes, cotisations sociales
\end{itemize}

\section{Les principaux marchés}

\subsection{Marché des biens et services}
\pointcle{Échange : Biens/services (flux réels) contre Paiements (flux monétaires)}
\begin{itemize}
    \item Entreprises → Ménages : produits et services
    \item Ménages → Entreprises : chiffre d'affaires
    \item Décisions influencées par la conjoncture économique
\end{itemize}

\subsection{Marché du travail}
\pointcle{Échange : Travail (flux réels) contre Salaires (flux monétaires)}
\begin{itemize}
    \item Ménages → Entreprises : force de travail
    \item Entreprises → Ménages : rémunération
    \item Rencontre offre et demande de travail
\end{itemize}

\subsection{Marché financier}
\pointcle{Lieu d'émission et d'échange des valeurs mobilières (actions, obligations)}
\begin{itemize}
    \item \textbf{Marché primaire} : émission de nouvelles valeurs
    \item \textbf{Marché secondaire} : échange de titres existants
    \item Financement externe direct (sans intermédiaire bancaire)
\end{itemize}

\section{Relations entreprise-État}

\pointcle{L'État fournit services non marchands, prestations sociales et subventions}

\subsection{Ce que l'État apporte aux entreprises}
\begin{itemize}
    \item Services publics (infrastructures, éducation, sécurité)
    \item Subventions et aides publiques
    \item Cadre juridique et réglementaire
\end{itemize}

\exemple{Bpifrance : banque publique qui finance et accompagne les projets d'entreprises (innovation, international, fonds propres)}

\subsection{Ce que les entreprises apportent à l'État}
\begin{itemize}
    \item Impôt sur les sociétés
    \item Cotisations sociales
    \item Taxes diverses
    \item Création d'emplois et de richesse
\end{itemize}

\section{Points clés à retenir}

\pointcle{L'entreprise est au centre d'un réseau de relations avec tous les agents économiques}

\pointcle{Ces relations créent des flux réels (biens, services, travail) et des flux monétaires (paiements, salaires)}

\pointcle{La conjoncture économique influence les comportements de consommation et d'investissement}

\pointcle{Le numérique transforme les relations client et nécessite de nouvelles approches marketing}

\pointcle{Les entreprises collaborent avec leurs fournisseurs pour améliorer performance et création de valeur}

\section{Mots-clés à connaître}
\begin{itemize}
    \item Agents économiques
    \item Production marchande/non marchande  
    \item Circuit économique
    \item Flux réels/monétaires
    \item Marché des biens et services
    \item Marché du travail
    \item Marché des capitaux
    \item Marché financier
    \item Conjoncture économique
    \item Financement externe direct/indirect
\end{itemize}

\end{document}